\section{When judges are participants}
\label{sec:boundary-cases}

The evaluator role has boundaries. If the study asks how linguists respond, how expertise changes judgments, how quickly people process a contrast, how dialect background affects acceptability, or how training changes ratings, then the judges belong to the target population. In such cases, the design studies their responses as responses, so they function as participants or participant-like contributors.

The clearest limiting case sits at the other end. A single-authored syntax article rests on the author's own judgments: the author builds the examples and marks them, and the marks carry the argument. If supplying a judgment relevant to the research question made one a participant, the author would be a participant in their own study, owing self-consent, a right to withdraw, and review of the author as a subject. No plausible ethics regime treats this as ordinary participant research, and on that reading the field's pre-experimental history counts as unconsented self-experimentation. Producing the judgment can't be the trigger.

Parity shows what the trigger is instead. Let the author judge fifty sentences and a colleague judge the same fifty; the judgments do identical evidential work, and only the producer differs. The line falls on the focus of inquiry, not the human act. A literalist might exempt the author for free, since no one else is studied, while counting the colleague a minimal participant; but that relocates the question rather than answering it, because production is constant across the two cases and the one thing that varies, who produced the verdict, has nothing to do with what the study is about. This doesn't make a lone intuition good evidence. That it can be unreliable is the measurement worry of \autoref{sec:problem}, met by calibration and replication, not by recasting the author as a research subject.

A single project can hold both roles. Experts may screen or validate materials as evaluators, while naive speakers then rate those materials as participants. If the study goes on to analyze the experts' own distribution of responses, their expertise effects, or their disagreement with naive speakers as a finding, the experts have become participants for that part of the design.

Because the role follows the design, it can look manipulable: relabel a judgment study as \enquote{evaluation} and the protections fall away. What blocks the move is that the role is fixed by the target of the research question, not by the label a researcher prefers. Take the hard case: thirty trained linguists rate two hundred sentences and report the reliability. If the ratings certify the status of the items, the design is evaluation. If the distribution of ratings is itself the finding, the linguists are the object of study, and they're participants. The diagnostic is what the judgment is evidence of, settled before any label is applied.

The distinction cuts both ways, against overreach and against special pleading. Expertise doesn't exempt a study from human-participant review when the human is the object of inquiry. But neither does the mere presence of human expertise convert every scholarly evaluation into participant research.
